\documentclass[aspectratio=169,11pt]{beamer}

\usetheme{Madrid}
\usecolortheme{default}
\usefonttheme{professionalfonts}
\setbeamertemplate{navigation symbols}{}
\setbeamertemplate{footline}[frame number]

\usepackage{amsmath, amssymb, booktabs}

\title{Professional Labor Markets And Public-Sector Design}
\subtitle{Assignment Rules, Career Ladders, And Public Objectives}
\author{Special Topic 7: Labor Market Design, Contracting, and Mechanism Design -- Week 5}
\date{}

\begin{document}

\begin{frame}
  \titlepage
\end{frame}

\begin{frame}{Central question}
\begin{itemize}
    \item What can structured professional and public-sector labor markets teach us about mechanism design under staffing, fairness, transparency, and career constraints?
    \item Week 5 keeps the labor object visible: workers are assigned to jobs, training paths, units, schools, and promotion ladders.
    \item The design problem is not only stability or strategy-proofness.
    \item It is rule $\rightarrow$ labor margin $\rightarrow$ public objective $\rightarrow$ worker welfare $\rightarrow$ dynamic career response.
\end{itemize}
\end{frame}

\begin{frame}{Professional and public-sector design principles}
\small
\begin{tabular}{p{0.26\linewidth} p{0.31\linewidth} p{0.31\linewidth}}
\toprule
Institutional feature & Labor margin & Design implication \\
\midrule
Compressed wages & nonprice allocation & assignment and priorities matter more \\
Public staffing targets & location/unit fill & efficiency is not only preference rank \\
Transparency requirements & legitimacy and fairness & rules must be explainable \\
Service obligations & entry, retention, exit & contracts change participation \\
Training pipelines & career value & initial matches are dynamic \\
Internal ladders & mobility, promotion & reassignment rules matter after entry \\
\bottomrule
\end{tabular}
\end{frame}

\begin{frame}{Rule vector and welfare objects}
\small
\[
  \rho = (A, P, C, L)
\]
\begin{itemize}
    \item $A$: assignment, clearing, rotation, or transfer rule.
    \item $P$: priority system, eligibility, reserves, scoring, staffing weights.
    \item $C$: contract terms, service obligations, probation, wage rigidity.
    \item $L$: internal ladder, promotion, transfer, training, on-the-job search.
\end{itemize}
\[
W_i(\rho)=u_i(j_i)+\tau_i(j_i,\rho)+g_i(L,\rho)
-c_i(j_i,\rho)-\lambda_i r_i(\rho)+b_i^{out}(\rho)
\]
\begin{itemize}
    \item Worker welfare includes assignment, training, mobility, risk, and outside options.
\end{itemize}
\end{frame}

\begin{frame}{Public objective and institutional detail}
\small
\[
  S(\rho)=
  \text{match quality}
  + \text{priority-site staffing}
  + \text{public-service output}
  + \text{retention}
  - \text{unfairness or opacity costs}
\]
\begin{itemize}
    \item The same matching rule can behave differently across residency, schools, civil service, and military labor markets.
    \item Wage rigidity makes priorities and assignment rules substitute for prices.
    \item Service obligations change outside options and continuation value.
    \item Staffing targets change what ``efficient'' means.
    \item Training and promotion systems make initial assignment dynamically consequential.
\end{itemize}
\end{frame}

\begin{frame}{Teacher and public-service assignment}
\begin{itemize}
    \item Teacher placement is a labor-market design problem, not only a school-choice analogy.
    \item The rule affects worker location choices, match quality, retention, training value, and public-service output.
    \item Public-service settings require explicit tradeoffs between teacher preferences, school priorities, geographic staffing, transparency, and student access.
    \item Recruitment contracts affect who enters; assignment rules affect where they work; career ladders affect whether they stay.
    \item Empirical design must separate worker preference satisfaction from service delivery and staffing equity.
\end{itemize}
\end{frame}

\begin{frame}{Internal ladders, training, and on-the-job search}
\small
\begin{tabular}{p{0.27\linewidth} p{0.33\linewidth} p{0.28\linewidth}}
\toprule
Dynamic margin & Observable object & Design question \\
\midrule
Training & credentials, completion, specialization & did assignment build human capital? \\
Promotion & rank, wage grade, leadership role & did the rule change career paths? \\
Transfer & internal moves, bids, rotations & can mismatch be repaired? \\
On-the-job search & applications, bids, exit threats & does the ladder retain workers? \\
Managerial discretion & rankings, overrides, evaluations & information or favoritism? \\
\bottomrule
\end{tabular}
\vspace{0.2cm}
\begin{itemize}
    \item Initial match quality is incomplete if later mobility, promotion, and training value change.
\end{itemize}
\end{frame}

\begin{frame}{Army and military assignment design}
\begin{itemize}
    \item Military assignment combines preferences, priorities, contracts, branch capacity, service obligations, and career pipelines.
    \item Branch-of-choice contracts show why matching with contracts is a labor-market object.
    \item Deferred acceptance can improve preference alignment, but implementation depends on branch priorities, manager behavior, and strategic communication.
    \item Assignment affects training, specialization, promotion, retention, and later internal mobility.
    \item The welfare claim must distinguish officer welfare, unit staffing, organizational mission, and long-run career outcomes.
\end{itemize}
\end{frame}

\begin{frame}{In and out of public employment}
\begin{itemize}
    \item Public-sector mechanisms operate only through workers who apply, accept, stay, transfer, or exit.
    \item Wage scales, pensions, security, mission value, probation, and promotion rules shape selection into the applicant pool.
    \item Service obligations and internal mobility affect whether workers remain after mismatch.
    \item Private-sector options affect sorting, wage pressure, and the interpretation of retention.
    \item Entry and exit are mechanism-design margins, not background attrition.
\end{itemize}
\end{frame}

\begin{frame}{Methods and data}
\small
\begin{tabular}{p{0.31\linewidth} p{0.29\linewidth} p{0.28\linewidth}}
\toprule
Design & Data & Identifies \\
\midrule
Mechanism randomization & market-by-market assignments & effect of rule change \\
Administrative redesign & preferences, assignments, outcomes & before/after rule differences \\
Ranking or score cutoff & applicant scores and later outcomes & local effect of entry or priority \\
Recruitment randomization & applicants plus performance & selection versus effort \\
Internal market data & bids, transfers, promotions & career mobility and mismatch repair \\
Surveys plus admin data & preferences and realized moves & objective alignment \\
\bottomrule
\end{tabular}
\vspace{0.2cm}
\begin{itemize}
    \item Best papers name the rule, labor margin, counterfactual, welfare object, public objective, and main threat.
\end{itemize}
\end{frame}

\begin{frame}{Research Lab design}
\begin{itemize}
    \item Primary anchor: deferred acceptance in Army officer labor markets.
    \item Challenge anchor: matching Teach For America teachers to schools.
    \item Reproduce: compare deferred-acceptance-style and manager-directed assignment in synthetic structured labor-market data.
    \item Diagnose: separate preference rank, staffing, priority alignment, training value, retention, internal mobility, and outside options.
    \item Transfer: adapt the design to teacher placement, public-service recruitment, civil-service entry, internal talent markets, or public-sector exit.
    \item Deliverable: rule, labor margin, counterfactual, data requirement, welfare object, public objective, and main threat.
\end{itemize}
\end{frame}

\begin{frame}{Takeaways}
\begin{itemize}
    \item Professional and public-sector labor markets make design choices visible.
    \item Assignment rules interact with wage rigidity, service obligations, training, internal ladders, and outside options.
    \item Military and teacher labor markets are empirical laboratories for staffing, fairness, transparency, and worker welfare.
    \item A credible welfare claim must track both public objectives and worker career outcomes.
    \item Week 6 turns these ingredients into full empirical research designs.
\end{itemize}
\end{frame}

\end{document}
